\documentclass[UTF8,12pt]{ctexart}
\usepackage[a4paper,margin=24mm]{geometry}
\usepackage{amsmath,amssymb,bm,graphicx,adjustbox}
\newcommand{\fitmath}[1]{\begin{adjustbox}{max width=\linewidth}$\displaystyle #1$\end{adjustbox}}
\setlength{\parindent}{0pt}
\setlength{\parskip}{.7em}
\sloppy
\allowdisplaybreaks
\begin{document}
\section*{1993 年考研数学三 · 真题}
\subsection*{1993 年全国硕士研究生招生考试试题（试卷 IV）}

\subsection*{一、填空题（本题共 5 小题，每小题 3 分，满分 15 分）}

（1）\(\displaystyle \lim_{x\to\infty}\dfrac{\displaystyle 3x^2+5}{\displaystyle 5x+3}\sin\dfrac{\displaystyle 2}{\displaystyle x}=\)\_\_\_\_\_\_。


（2）已知 \(\displaystyle y=f\left(\dfrac{\displaystyle 3x-2}{\displaystyle 3x+2}\right)\)，\(\displaystyle f^{\prime}(x)=\arctan x^2\)，则 \(\displaystyle \left.\dfrac{\displaystyle dy}{\displaystyle dx}\right|_{x=0}=\)\_\_\_\_\_\_。


（3）级数 \(\displaystyle \sum_{n=0}^{\infty}\dfrac{\displaystyle (\ln3)^n}{\displaystyle 2^n}\) 的和为\_\_\_\_\_\_。


（4）设 4 阶方阵 \(\displaystyle A\) 的秩为 2，则其伴随矩阵 \(\displaystyle A^*\) 的秩为\_\_\_\_\_\_。


（5）（超纲题）设总体 \(\displaystyle X\) 的方差为 1，根据来自 \(\displaystyle X\) 的容量为 100 的简单随机样本，测得样本均值为 5，则 \(\displaystyle X\) 的数学期望的置信度近似等于 0.95 的置信区间为\_\_\_\_\_\_。


\subsection*{二、选择题（本题共 5 小题，每小题 3 分，满分 15 分）}

（1）设函数 \(\displaystyle f(x)=\begin{cases}\sqrt{|x|}\sin\dfrac{\displaystyle 1}{\displaystyle x^2},\allowbreak &x\ne0,\allowbreak \\0,\allowbreak &x=0,\allowbreak \end{cases}\) 则 \(\displaystyle f(x)\) 在 \(\displaystyle x=0\) 处（　）。


（A）极限不存在。 （B）极限存在但不连续。 （C）连续但不可导。 （D）可导。


（2）设 \(\displaystyle f(x)\) 为连续函数，且 \(\displaystyle F(x)=\displaystyle \int_{1/x}^{\ln x}f(t)\,dt\)，则 \(\displaystyle F^{\prime}(x)\) 等于（　）。


（A）\(\displaystyle \dfrac{\displaystyle 1}{\displaystyle x} f(\ln x)+\dfrac{\displaystyle 1}{\displaystyle x^2}f\left(\dfrac{\displaystyle 1}{\displaystyle x}\right)\)。 （B）\(\displaystyle f(\ln x)+f\left(\dfrac{\displaystyle 1}{\displaystyle x}\right)\)。


（C）\(\displaystyle \dfrac{\displaystyle 1}{\displaystyle x} f(\ln x)-\dfrac{\displaystyle 1}{\displaystyle x^2}f\left(\dfrac{\displaystyle 1}{\displaystyle x}\right)\)。 （D）\(\displaystyle f(\ln x)-f\left(\dfrac{\displaystyle 1}{\displaystyle x}\right)\)。


（3）n 阶方阵 A 具有 n 个不同的特征值是 A 与对角阵相似的（　）。


（A）充分必要条件。 （B）充分而非必要条件。 （C）必要而非充分条件。 （D）既非充分也非必要条件。


（4）（原考题有误，见参考答案说明）设两事件 \(\displaystyle A\) 与 \(\displaystyle B\) 满足 \(\displaystyle P(B\mid A)=1\)，则（　）。


（A）\(\displaystyle A\) 是必然事件。 （B）\(\displaystyle P(B\mid\bar A)=0\)。 （C）\(\displaystyle A\supset B\)。 （D）\(\displaystyle A\subset B\)。


（5）设随机变量 \(\displaystyle X\) 的概率密度为 \(\displaystyle \varphi(x)\)，且 \(\displaystyle \varphi(-x)=\varphi(x)\)，\(\displaystyle F(x)\) 为 \(\displaystyle X\) 的分布函数，则对任意实数 \(\displaystyle a\)，有（　）。


（A）\(\displaystyle F(-a)=1-\displaystyle \int_0^a\varphi(x)\,dx\)。 （B）\(\displaystyle F(-a)=\dfrac{\displaystyle 1}{\displaystyle 2}-\displaystyle \int_0^a\varphi(x)\,dx\)。 （C）\(\displaystyle F(-a)=F(a)\)。 （D）\(\displaystyle F(-a)=2F(a)-1\)。


\subsection*{试卷 IV（续）}

三、（本题满分 5 分）设 \(\displaystyle z=f(x,\allowbreak y)\) 是由方程 \(\displaystyle z-y-x+xe^{z-y-x}=0\) 所确定的二元函数，求 \(\displaystyle dz\)。


四、（本题满分 7 分）已知 \(\displaystyle \lim_{x\to\infty}\left(\dfrac{\displaystyle x-a}{\displaystyle x+a}\right)^x=\displaystyle \int_a^{+\infty}4x^2e^{-2x}\,dx\)，求常数 \(\displaystyle a\) 的值。


五、（本题满分 9 分）设某产品的成本函数为 \(\displaystyle C=aq^2+bq+c\)，需求函数为 \(\displaystyle q=\dfrac{\displaystyle 1}{\displaystyle e}(d-p)\)，其中 \(\displaystyle C\) 为成本，\(\displaystyle q\) 为需求量（即产量），\(\displaystyle p\) 为单价，\(\displaystyle a,\allowbreak b,\allowbreak c,\allowbreak d,\allowbreak e\) 都是正的常数，且 \(\displaystyle d>b\)。求：（1）利润最大时的产量及最大利润；（2）需求对价格的弹性；（3）需求对价格弹性的绝对值为 1 时的产量。


六、（本题满分 8 分）假设：（1）函数 \(\displaystyle y=f(x)\)（\(\displaystyle 0\leq x<+\infty\)）满足条件 \(\displaystyle f(0)=0\) 和 \(\displaystyle 0\leq f(x)\leq e^x-1\)；（2）平行于 \(\displaystyle y\) 轴的动直线 \(\displaystyle MN\) 与曲线 \(\displaystyle y=f(x)\) 和 \(\displaystyle y=e^x-1\) 分别相交于点 \(\displaystyle P_1\) 和 \(\displaystyle P_2\)；（3）曲线 \(\displaystyle y=f(x)\)，直线 \(\displaystyle MN\) 与 \(\displaystyle x\) 轴所围封闭图形的面积 \(\displaystyle S\) 恒等于线段 \(\displaystyle P_1P_2\) 的长度。求函数 \(\displaystyle y=f(x)\) 的表达式。


七、（本题满分 6 分）假设函数 \(\displaystyle f(x)\) 在 \(\displaystyle [0,\allowbreak 1]\) 上连续，在 \(\displaystyle (0,\allowbreak 1)\) 内二阶可导，过点 \(\displaystyle A(0,\allowbreak f(0))\) 与 \(\displaystyle B(1,\allowbreak f(1))\) 的直线与曲线 \(\displaystyle y=f(x)\) 相交于点 \(\displaystyle C(c,\allowbreak f(c))\)，其中 \(\displaystyle 0<c<1\)。证明：在 \(\displaystyle (0,\allowbreak 1)\) 内至少存在一点 \(\displaystyle \xi\)，使 \(\displaystyle f^{\prime\prime}(\xi)=0\)。


八、（本题满分 10 分）\(\displaystyle k\) 为何值时，线性方程组 \(\displaystyle \begin{cases}x_1+x_2+kx_3=4,\allowbreak \\-x_1+kx_2+x_3=k^2,\allowbreak \\x_1-x_2+2x_3=-4\end{cases}\) 有唯一解、无解、有无穷多解？在有解情况下，求出其全部解。


九、（本题满分 9 分）设二次型 \(\displaystyle f=x_1^2+x_2^2+x_3^2+2\alpha x_1x_2+2\beta x_2x_3+2x_1x_3\) 经正交变换 \(\displaystyle x=Py\) 化成 \(\displaystyle f=y_2^2+2y_3^2\)，其中 \(\displaystyle x=(x_1,\allowbreak x_2,\allowbreak x_3)^{\mathrm T}\) 和 \(\displaystyle y=(y_1,\allowbreak y_2,\allowbreak y_3)^{\mathrm T}\) 是 3 维列向量，\(\displaystyle P\) 是 3 阶正交矩阵。试求常数 \(\displaystyle \alpha,\allowbreak \beta\)。


十、（本题满分 8 分）设随机变量 \(\displaystyle X\) 和 \(\displaystyle Y\) 同分布，\(\displaystyle X\) 的概率密度为 \(\displaystyle f(x)=\begin{cases}\dfrac{\displaystyle 3}{\displaystyle 8}x^2,\allowbreak &0<x<2,\allowbreak \\0,\allowbreak &\text{其他}.\end{cases}\)


\subsection*{试卷 IV（续）}

十、（续）（1）已知事件 \(\displaystyle A=\{X>a\}\) 和 \(\displaystyle B=\{Y>a\}\) 独立，且 \(\displaystyle P(A\cup B)=\dfrac{\displaystyle 3}{\displaystyle 4}\)，求常数 \(\displaystyle a\)；（2）求 \(\displaystyle \dfrac{\displaystyle 1}{\displaystyle X^2}\) 的数学期望。


十一、（本题满分 8 分）假设一大型设备在任何长为 \(\displaystyle t\) 的时间内发生故障的次数 \(\displaystyle N(t)\) 服从参数为 \(\displaystyle \lambda t\) 的泊松分布，（1）求相继两次故障之间时间间隔 \(\displaystyle T\) 的概率分布；（2）求在设备已经无故障工作 8 小时的情形下，再无故障运行 8 小时的概率 \(\displaystyle Q\)。


\subsection*{1993 年全国硕士研究生招生考试试题（试卷 V）}

\subsection*{一、填空题（本题共 5 小题，每小题 3 分，满分 15 分）}

（1）\(\displaystyle \lim_{n\to\infty}\left[\sqrt{1+2+\cdots+n}-\sqrt{1+2+\cdots+(n-1)}\right]=\)\_\_\_\_\_\_。


（2）已知 \(\displaystyle y=f\left(\dfrac{\displaystyle 3x-2}{\displaystyle 3x+2}\right)\)，\(\displaystyle f^{\prime}(x)=\arcsin x^2\)，则 \(\displaystyle \left.\dfrac{\displaystyle dy}{\displaystyle dx}\right|_{x=0}=\)\_\_\_\_\_\_。


（3）\(\displaystyle \int\dfrac{\displaystyle dx}{\displaystyle (2-x)\sqrt{1-x}}=\)\_\_\_\_\_\_。


（4）【同试卷 IV 第一、（4）题】


（5）设 10 件产品有 4 件不合格品，从中任取两件，已知所取两件产品中有一件是不合格品，则另一件也是不合格品的概率为\_\_\_\_\_\_。


\subsection*{二、选择题（本题共 5 小题，每小题 3 分，满分 15 分）}

（1）【同试卷 IV 第二、（1）题】


（2）【同试卷 IV 第二、（2）题】


（3）若 \(\displaystyle \alpha_1,\allowbreak \alpha_2,\allowbreak \alpha_3,\allowbreak \beta_1,\allowbreak \beta_2\) 都是 4 维列向量，且 4 阶行列式 \(\displaystyle |\alpha_1,\allowbreak \alpha_2,\allowbreak \alpha_3,\allowbreak \beta_1|=m,\allowbreak \ |\alpha_1,\allowbreak \alpha_2,\allowbreak \beta_2,\allowbreak \alpha_3|=n\)，则 4 阶行列式 \(\displaystyle |\alpha_3,\allowbreak \alpha_2,\allowbreak \alpha_1,\allowbreak \beta_1+\beta_2|\) 等于（　）。


（A）\(\displaystyle m+n\)。 （B）\(\displaystyle -(m+n)\)。 （C）\(\displaystyle n-m\)。 （D）\(\displaystyle m-n\)。


（4）设 \(\displaystyle \lambda=2\) 是非奇异矩阵 \(\displaystyle A\) 的一个特征值，则矩阵 \(\displaystyle \left(\dfrac{\displaystyle 1}{\displaystyle 3} A^2\right)^{-1}\) 有一特征值等于（　）。


（A）\(\displaystyle \dfrac{\displaystyle 4}{\displaystyle 3}\)。 （B）\(\displaystyle \dfrac{\displaystyle 3}{\displaystyle 4}\)。 （C）\(\displaystyle \dfrac{\displaystyle 1}{\displaystyle 2}\)。 （D）\(\displaystyle \dfrac{\displaystyle 1}{\displaystyle 4}\)。


（5）设随机变量 \(\displaystyle X\) 与 \(\displaystyle Y\) 均服从正态分布，\(\displaystyle X\sim N(\mu,\allowbreak 4^2),\allowbreak \ Y\sim N(\mu,\allowbreak 5^2)\)，记 \(\displaystyle p_1=P\{X\leq\mu-4\},\allowbreak \ p_2=P\{Y\geq\mu+5\}\)，则（　）。


（A）对任何实数 \(\displaystyle \mu\)，都有 \(\displaystyle p_1=p_2\)。 （B）对任何实数 \(\displaystyle \mu\)，都有 \(\displaystyle p_1<p_2\)。 （C）只对 \(\displaystyle \mu\) 的个别值，才有 \(\displaystyle p_1=p_2\)。 （D）对任何实数 \(\displaystyle \mu\)，都有 \(\displaystyle p_1>p_2\)。


三、（本题满分 5 分）【同试卷 IV 第三题】


四、（本题满分 7 分）【同试卷 IV 第四题】


\subsection*{试卷 V（续）}

五、（本题满分 7 分）已知某厂生产 \(\displaystyle x\) 件产品的成本为 \(\displaystyle C=25000+200x+\dfrac{\displaystyle 1}{\displaystyle 40}x^2\)（元），问：（1）要使平均成本最小，应生产多少件产品？（2）若产品以每件 500 元售出，要使利润最大，应生产多少件产品？


六、（本题满分 6 分）设 \(\displaystyle p,\allowbreak q\) 是大于 1 的常数，并且 \(\displaystyle \dfrac{\displaystyle 1}{\displaystyle p}+\dfrac{\displaystyle 1}{\displaystyle q}=1\)。证明：对于任意的 \(\displaystyle x>0\)，有 \(\displaystyle \dfrac{\displaystyle 1}{\displaystyle p} x^p+\dfrac{\displaystyle 1}{\displaystyle q}\geq x\)。


七、（本题满分 13 分）运用导数的知识作函数 \(\displaystyle y=(x+6)e^{1/x}\) 的图形。


八、（本题满分 8 分）已知 3 阶矩阵 \(\displaystyle A\) 的逆矩阵为 \(\displaystyle A^{-1}=\begin{pmatrix}1&1&1\\1&2&1\\1&1&3\end{pmatrix}\)，试求其伴随矩阵 \(\displaystyle A^*\) 的逆矩阵。


九、（本题满分 8 分）设 \(\displaystyle A\) 是 \(\displaystyle m\times n\) 矩阵，\(\displaystyle B\) 是 \(\displaystyle n\times m\) 矩阵，\(\displaystyle E\) 是 \(\displaystyle n\) 阶单位矩阵（\(\displaystyle m>n\)），已知 \(\displaystyle BA=E\)。试判断 \(\displaystyle A\) 的列向量组是否线性相关？为什么？


十、（本题满分 8 分）设随机变量 \(\displaystyle X\) 和 \(\displaystyle Y\) 独立，都在区间 \(\displaystyle [1,\allowbreak 3]\) 上服从均匀分布。引进事件 \(\displaystyle A=\{X\leq a\},\allowbreak \ B=\{Y>a\}\)。（1）已知 \(\displaystyle P(A\cup B)=\dfrac{\displaystyle 7}{\displaystyle 9}\)，求常数 \(\displaystyle a\)；（2）求 \(\displaystyle \dfrac{\displaystyle 1}{\displaystyle X}\) 的数学期望。


十一、（本题满分 8 分）【同试卷 IV 第十一题】

\end{document}
